\documentclass[11pt]{article}

\usepackage[margin=1in]{geometry}
\usepackage[utf8]{inputenc}
\usepackage[T1]{fontenc}
\usepackage{hyperref}
\usepackage{url}
\usepackage{booktabs}
\usepackage{amsfonts}
\usepackage{amsmath}
\usepackage{microtype}
\usepackage{enumitem}
\usepackage{xcolor}
\usepackage[numbers,sort&compress]{natbib}
\usepackage{array}

\hypersetup{
  colorlinks=true,
  linkcolor=blue!60!black,
  citecolor=blue!60!black,
  urlcolor=blue!60!black
}

% ── Compact layout ──────────────────────────────────────────────
\setlength{\parskip}{4pt plus 1pt minus 1pt}
\setlength{\parindent}{0pt}

\title{%
  \textbf{Pre-Registration:} Interaction Effects, Scaling Laws, and\\
  Pareto-Optimal Profiles in Retrieval-Augmented\\
  Generation Pipelines%
}

\author{%
  [Authors anonymized for review]%
}

\date{}

\begin{document}
\maketitle
\thispagestyle{empty}

\vspace{-1.5em}
\begin{center}
\begin{tabular}{@{}rl@{}}
  \textbf{Date of pre-registration:} & [DATE] \\
  \textbf{Design matrix SHA-256:}    & \texttt{31a0dc197b46...df1728d6} \\
  \textbf{Software:}                 & Python 3.11, NumPy 2.3.5, SciPy 1.16.3, seed\,$=$\,42 \\
\end{tabular}
\end{center}

% ════════════════════════════════════════════════════════════════
%  Abstract
% ════════════════════════════════════════════════════════════════

\begin{abstract}
\noindent
\textbf{In one sentence:} we commit, before collecting data, to
testing whether a practitioner can build a good Retrieval-Augmented
Generation (RAG) pipeline by tuning components one at a time on a
small corpus, the way the literature implicitly assumes.
We answer this through three pre-registered hypotheses.
\textbf{H1 (additivity)} asks whether the eight pipeline components
combine additively, or whether their interactions are large enough
that single-component ablations cannot be trusted to compose.
\textbf{H2 (a balanced recipe)} asks whether one fixed
\textsc{balanced} configuration sweeps a wider region of the
(faithfulness, latency, cost) tradeoff envelope than a
quality-only configuration, or whether each deployment must
hand-tune for its own operating point.
\textbf{H3 (scaling)} asks whether the small-corpus profile ranking
that drives most published RAG ablations survives a $40\times$
corpus scale-up, or whether the small-corpus winner is not the
production-scale winner.
The pre-registered study commits to a 242-configuration D-optimal
fractional factorial design spanning 8 pipeline dimensions, 5
benchmark datasets (2{,}300 queries per configuration), and a
three-track evaluation protocol covering retrieval, end-to-end
generation, and hallucination resistance.  To our knowledge, this is
the first pre-registered factorial evaluation of RAG pipeline
composition---$4.6\times$ the scale of the largest prior controlled
RAG study~\citep{jin2024flashrag}.  We publish this document before
experimental data collection to establish hypothesis commitment and
invite community feedback on model selection, benchmark composition,
and statistical methodology.
\end{abstract}

% ════════════════════════════════════════════════════════════════
\section{Motivation and Rationale}
% ════════════════════════════════════════════════════════════════

Retrieval-Augmented Generation~\citep{lewis2020rag} has become the
dominant paradigm for grounding large language models in external
knowledge.  A rich toolkit of individually validated
techniques---dense retrieval, hybrid sparse--dense fusion, cross-encoder
reranking, HyDE~\citep{gao2023hyde}, CRAG~\citep{yan2024crag},
Self-RAG~\citep{asai2023selfrag}---is available to practitioners, yet
each technique is evaluated in isolation against a simple baseline.
Practitioners who combine several techniques simultaneously have no
empirical basis for predicting whether their gains add, cancel, or
amplify.

This is not merely a reporting gap.
\citet{wang2024searching} observe that end-to-end RAG performance
depends on coordinated pipeline design rather than any single component;
\citet{chen2025mmoarag} confirm the presence of interaction effects by
showing that jointly optimized RAG components outperform independently
tuned ones.  Yet neither work measures the sign, magnitude, or
domain-dependence of these interactions under controlled factorial
conditions---precisely the question a practitioner faces when choosing a
configuration for production.

Pre-registration prevents post-hoc hypothesis mining, a known problem in
ML benchmarking~\citep{dodge2019importance,bouthillier2021accounting}.
By committing to hypotheses, analysis plans, and exclusion criteria
before data collection, we ensure that our confirmatory findings are
distinguishable from exploratory ones.  We additionally publish this
document to invite community input on experimental choices that remain
open (Section~\ref{sec:models}).

% ════════════════════════════════════════════════════════════════
\section{Pre-Registered Hypotheses}
\label{sec:hypotheses}
% ════════════════════════════════════════════════════════════════

We pre-register three hypotheses that together ask whether RAG
pipelines can be designed compositionally, picking up the gap
identified in Section~1.  \textbf{H1} asks the
\emph{structural} question: do the eight pipeline components combine
additively, or do they interact non-trivially?  If they interact, the
implicit assumption behind every single-component ablation in the
literature fails.  \textbf{H2} asks the \emph{design} question:
assuming components do interact, can a single fixed configuration
still serve practitioners across the (faithfulness, latency, cost)
tradeoff space, or must each deployment hand-tune for its own operating
point?  \textbf{H3} asks the \emph{scaling} question: do the answers
to H1 and H2 survive a $40\times$ scale-up of the corpus, or are
small-corpus best practices misleading at production scale?  The three
together determine whether a fixed ``best-practice RAG pipeline'' is
even a coherent target.

A summary of the three before the formal statements:

\begin{table}[h]
\centering
\footnotesize
\begin{tabular}{@{}p{0.8cm}p{4.4cm}p{4.4cm}p{4.0cm}@{}}
\toprule
& Research question & What confirming rejects & Pre-registered test \\
\midrule
\textbf{H1} & Do components combine additively? & Single-component ablation as evidence & $\eta^2 \geq 0.20$, Type-III ANOVA F-test \\
\addlinespace
\textbf{H2} & Does a balanced configuration sweep a wider tradeoff envelope? & ``You must pick one axis to optimize'' & Hypervolume, Wilcoxon signed-rank \\
\addlinespace
\textbf{H3} & Do RAG profile rankings flip with corpus scale? & Small-corpus ablations predict production deployment & Per-profile power-law slopes, bootstrap CIs, Kendall's $\tau$ \\
\bottomrule
\end{tabular}
\end{table}

\paragraph{H1 --- Components do not combine additively.}

\textit{Research question.} Do the eight pipeline components combine
additively, or do non-trivial interactions arise that violate the
additivity assumption behind isolated-component ablations?

\textit{Pre-registered claim.} Two-way interactions between the eight
factorial dimensions explain $\geq 20\%$ of the variance in
faithfulness (RAGAS, 0--1 scale) after main effects are removed.
Variance is reported as \emph{partial $\eta^2$} --- the share of
residual variance (after main effects) attributed to interaction
terms --- in a \emph{Type-III ANOVA}, an ANOVA decomposition that
estimates each term's contribution after partialling out all others
and is robust to unbalanced designs.  The 20\% bar is Cohen's ``large
effect'' threshold for a single interaction term~\citep{cohen1988statistical}.

\textit{Pre-registered test.} F-test on the joint set of all 28
two-way interaction terms, $\alpha = 0.05$.  Post-hoc Bonferroni
($\alpha = 0.05 / 28 = 0.00179$) per interaction pair.

\textit{What confirming and rejecting mean.} \emph{Confirming~H1}
falsifies the additivity assumption behind most single-component RAG
ablations in the literature: a designer cannot tune one component in
isolation and trust the gain to compose with another.
\emph{Rejecting~H1} (interaction terms contribute $<$5\% of variance)
means RAG components are approximately additive within the precision
of our 242-config sweep, and isolated ablations remain a defensible
practice.

\paragraph{H2 --- A balanced operating point exists.}

\textit{Research question.} Given that components interact, can a
single fixed configuration sweep a wider region of the (faithfulness,
latency, cost) tradeoff space than a quality-only configuration, or
must practitioners give up an axis to gain another?

\textit{Pre-registered claim.} The \textsc{balanced} operational
profile exhibits a strictly higher \emph{hypervolume indicator} (HV)
than the \textsc{max\_reliable} profile on $\geq 3$ of the 5 benchmark
datasets.  Hypervolume is the multi-objective generalization of
area-under-the-curve: the volume of the
$(\text{faithfulness}{\uparrow}, \text{latency}{\downarrow}, \text{cost}{\downarrow})$
region dominated by the profile's per-query tradeoff cloud, with
larger meaning a better envelope.

\textit{Pre-registered test.} Wilcoxon signed-rank --- a non-parametric
paired-sample test that does not assume normality --- on per-query
faithfulness scores (paired by query), Bonferroni-corrected across
the 5 datasets ($\alpha = 0.05 / 5 = 0.01$ per dataset).  HV is
computed with the exact 3D algorithm; objectives are normalized to
$[0,1]$ before the volume integration.

\textit{What confirming and rejecting mean.} \emph{Confirming~H2}
says a single fixed configuration can serve practitioners who do not
want to give up an axis: published profile recipes (e.g., \textsc{Balanced})
generalize across deployments.  \emph{Rejecting~H2}
(\textsc{max\_reliable} dominates in $\geq 4 / 5$ datasets) means
production RAG requires per-deployment tuning to its own tradeoff
preference, and our published profile recipes will not generalize.

\paragraph{H3 --- Profile rankings flip with corpus scale.}

\textit{Research question.} Does the best-performing RAG profile at
small corpus (where most published RAG ablations live) remain the
best-performing profile at production scale, or do profile rankings
change with corpus size?

\textit{Pre-registered claim.} Fitting NDCG@10 as a power law
$\text{NDCG@10}(N) = a \cdot N^{b}$ separately for each of the four
operational profiles (\textsc{Speed}, \textsc{Standard},
\textsc{Balanced}, \textsc{Max\_Reliable}), the cross-profile range
$b_{\max} - b_{\min}$ exceeds $0.05$ per decade of corpus growth
(i.e., the steepest- and shallowest-degrading profiles differ by at
least $5\%$ relative NDCG@10 per $10\times$ corpus increase),
analogous to the saturation behavior observed in neural scaling
laws~\citep{kaplan2020scaling}.  Equivalently, the profile-rank
ordering at the smallest tested corpus ($5$K) differs from the
ordering at $200$K on $\geq 1$ benchmark.

\textit{Pre-registered test.} Per-profile nonlinear least-squares
fit (\texttt{scipy.optimize.curve\_fit}); 95\% bootstrap CIs on each
$b$ from 10{,}000 resamples.  Profile-pair slope differences tested
via \emph{permutation test} (10{,}000 shuffles, $\alpha = 0.05$) ---
a non-parametric significance test that constructs the null
distribution by random label shuffling, valid without distributional
assumptions.  Profile-rank stability tested via Kendall's $\tau$
between the $5$K and $200$K profile orderings.

\textit{What confirming and rejecting mean.} \emph{Confirming~H3}
means small-corpus ablations cannot reliably predict production-scale
deployment choices: profile selection requires re-validation at the
deployment's target corpus size, and the field's reliance on
$\leq 50$K-corpus benchmarks is methodologically suspect for
deployment decisions.  \emph{Rejecting~H3} (slope range $\leq 0.02$
per decade, $\tau = 1$ on every benchmark) means small-corpus
benchmarks generalize, and the small-corpus best practice is also the
large-corpus best practice.

% ════════════════════════════════════════════════════════════════
\section{Experimental Design}
\label{sec:design}
% ════════════════════════════════════════════════════════════════

% ── Platform ──────────────────────────────────────────────────
\subsection{Platform}

\textsc{Lasagna} is a stage-isolated RAG evaluation platform: each
pipeline stage is implemented behind an abstract base class, and
configuration changes to one stage do not affect the execution of any
other stage.  This architectural invariant guarantees that the output
of stage~$i$ is a deterministic function of the output of stage~$i{-}1$
and the configuration parameters~$\boldsymbol{\theta}_i$ alone,
eliminating the coupling confounds that invalidate factorial
interpretation in conventional RAG implementations.

% ── Candidate Models ──────────────────────────────────────────
\subsection{Candidate Models}
\label{sec:models}

All configurations share a single generator LLM (held constant) and a
single evaluation judge.  We list primary and alternative candidates
to invite community input before experiments begin.

\begin{table}[h]
\centering
\footnotesize
\caption{Model candidates.  The generator is held constant across all
  configurations; the judge evaluates all RAGAS metrics.  Embedding
  models are factorial dimension~D1.}
\label{tab:models}
\begin{tabular}{@{}p{2.0cm}p{3.4cm}p{3.8cm}p{4.0cm}@{}}
\toprule
\textbf{Role} & \textbf{Primary} & \textbf{Alternatives} & \textbf{Selection rationale} \\
\midrule
Generator    & GPT-4.1 (OpenAI)
             & Claude~3.7 Sonnet (Anthropic), Gemini~2.0 (Google)
             & Held constant; cost vs.\ capability \\
\addlinespace
Judge        & GPT-4o-mini (OpenAI)
             & Claude~3.5 Haiku (Anthropic), Llama-3-70B (Meta)
             & Determinism at $T{=}0$; cost for ${\sim}$1.7M calls \\
\addlinespace
Embedding    & \multicolumn{2}{l}{MiniLM-L6 (384d), BGE-M3 (1024d), E5-large (1024d), BGE-large (1024d)}
             & Span lightweight to SOTA multilingual \\
\addlinespace
Reranker     & ms-marco-MiniLM-L6-v2
             & BGE-reranker-v2
             & Standard BEIR cross-encoder baseline \\
\bottomrule
\end{tabular}
\end{table}

We welcome suggestions on model selection before experiments begin.
In particular, we are open to: (1)~alternative generator models that
would increase reproducibility (e.g., open-weight models),
(2)~additional judge models for cross-validation of faithfulness
scores, and (3)~embedding models that would expand the capacity range.

% ── Factorial Dimensions ──────────────────────────────────────
\subsection{Factorial Dimensions}

\begin{table}[h]
\centering
\footnotesize
\caption{The eight factorial dimensions.  Baseline levels (marked~$^*$)
  are the simplest, lowest-cost options.}
\label{tab:dimensions}
\begin{tabular}{@{}llcl@{}}
\toprule
Dim & Stage & $k$ & Levels \\
\midrule
D1 & Embedding      & 4 & MiniLM-L6$^*$, BGE-M3, E5-large, BGE-large \\
D2 & Chunking       & 4 & fixed$^*$, semantic, sentence-window, hierarchical \\
D3 & Enrichment     & 4 & none$^*$, contextual-retrieval, RAPTOR, GraphRAG \\
D4 & Retrieval      & 3 & dense$^*$, hybrid-BM25, hybrid-SPLADE \\
D5 & Query aug.     & 5 & none$^*$, HyDE, RAG-Fusion, decomposition, step-back \\
D6 & Reranking      & 4 & none$^*$, single-CE, ensemble, MMR \\
D7 & Compression    & 2 & none$^*$, LLMLingua \\
D8 & Agentic gate   & 4 & none$^*$, CRAG, Self-RAG, Speculative-RAG \\
\bottomrule
\end{tabular}
\end{table}

The full factorial comprises
$4 \times 4 \times 4 \times 3 \times 5 \times 4 \times 2 \times 4 =
30{,}720$ configurations.

% ── D-Optimal Design ─────────────────────────────────────────
\subsection{D-Optimal Design}

We reduce 30{,}720 candidates to \textbf{242 configurations} using a
\emph{D-optimal fractional factorial} --- a design that selects a
subset of rows from the full factorial to give the most precise
estimates of every main effect and pairwise interaction, formally by
maximizing $\det(X^{\top} X)$ for the model matrix $X$ --- generated
by Fedorov coordinate
exchange~\citep{fedorov1972theory,atkinson2007optimum}.  The model
matrix includes 1~intercept $+$ 22~main-effect $+$ 209~two-way
interaction columns $=$ 232 parameters, with 10 residual degrees of
freedom.  Effect coding ensures orthogonal main
effects~\citep{box2005statistics,montgomery2017design}.  All 232 model
terms are estimable (rank verification), and no two-way interaction
correlates with any representative three-way interaction at
$|r| > 0.30$.

\textbf{Design matrix SHA-256:}
\texttt{31a0dc197b46a90819fc9178ea06668444ad9617ff4c8a86f3d557a4df1728d6}
(seed\,$=$\,42).

% ── Benchmarks ────────────────────────────────────────────────
\subsection{Benchmark Datasets}

\begin{table}[h]
\centering
\footnotesize
\caption{Planned benchmark datasets.  Total evaluation queries per
  configuration: 2{,}300.}
\label{tab:benchmarks}
\begin{tabular}{@{}llrrl@{}}
\toprule
ID & Dataset & Corpus & Eval queries & Relevance \\
\midrule
DS1 & MS~MARCO~v1     & 8.8M passages       & 500 sampled & BM25 sparse (TREC) \\
DS2 & HotpotQA        & 5.2M paragraphs      & 500 sampled & Supporting facts \\
DS3 & BEIR/FiQA       & 57K posts            & 500 sampled & Crowdsourced \\
DS4 & BEIR/SciFact    & 5K abstracts         & 300 (all)   & Expert binary \\
DS5 & CUAD (legal)    & 510 contracts        & 500 sampled & Clause-level \\
\bottomrule
\end{tabular}
\end{table}

These five benchmarks span general web (DS1), multi-hop reasoning
(DS2), financial opinion QA (DS3), scientific claim verification
(DS4), and legal contract analysis (DS5), covering diversity in domain,
relevance granularity, and query
complexity~\citep{thakur2021beir}.

% ── Metrics ───────────────────────────────────────────────────
\subsection{Metrics}

\textit{Retrieval:} NDCG@10, MRR, Recall@100.
\textit{End-to-end:} faithfulness, answer relevance
(RAGAS~v0.4~\citep{es2023ragas}, judged by GPT-4o-mini at $T{=}0$),
exact match, token~F1, p50/p95 latency, \$/query.
\textit{Multi-objective:} \emph{hypervolume indicator} (HV) ---
the volume of the
$(\text{faithfulness}{\uparrow}, \text{latency}{\downarrow}, \text{cost}{\downarrow})$
region dominated by a profile's per-query tradeoff cloud, the
multi-objective generalization of area-under-the-curve --- in the
faithfulness--latency--cost space, with reference point
$(0, 10{,}000\,\text{ms}, \$1.00)$.

% ════════════════════════════════════════════════════════════════
\section{Statistical Analysis Plan}
\label{sec:analysis}
% ════════════════════════════════════════════════════════════════

\paragraph{H1 (additivity).}
\emph{Plain reading:} fit one regression that includes every
pipeline component plus all 28 of their pairwise interactions, then
ask how much of the leftover variance the interactions explain.
Model:
\begin{small}
\begin{verbatim}
faithfulness ~ D1+D2+D3+D4+D5+D6+D7+D8
             + D1:D2+D1:D3+...+D7:D8   (28 two-way terms)
\end{verbatim}
\end{small}
\vspace{-0.5em}
Test statistic: F for the joint set of all interaction terms,
$F(209,\, n_\text{resid})$.  Effect size: partial
$\eta^2 = \text{SS}_\text{int} / (\text{SS}_\text{int} +
\text{SS}_\text{resid})$; threshold $\geq 0.20$.  Post-hoc:
Bonferroni ($\alpha = 0.05/28 = 0.00179$) per interaction pair, so
that any single interaction we flag as ``large'' has been corrected
for the 28 simultaneous comparisons.

\paragraph{H2 (a balanced recipe).}
\emph{Plain reading:} for each query, compare the
\textsc{balanced} profile's three-objective tradeoff cloud against
the \textsc{max\_reliable} profile's, and check on how many of the
five benchmarks \textsc{balanced} sweeps a strictly larger
hypervolume.  Wilcoxon signed-rank test on per-query faithfulness
(\textsc{balanced} vs.\ \textsc{max\_reliable}, paired by query).
Bonferroni across 5~datasets ($\alpha = 0.01$ per dataset) so that
any single ``win'' has been corrected for testing on five datasets
in parallel.  HV computed with the exact 3D algorithm; objectives
normalized to $[0,1]$ before integration.

\paragraph{H3 (scaling).}
\emph{Plain reading:} fit a separate scaling curve for each of the
four profiles and check whether the steepest- and shallowest-
degrading profiles drift apart fast enough that the small-corpus
winner is no longer the large-corpus winner.  Per-profile nonlinear
least-squares fit of $\text{NDCG@10}(N) = a \cdot N^b$ for each of
\textsc{Speed}, \textsc{Standard}, \textsc{Balanced},
\textsc{Max\_Reliable} using \texttt{scipy.optimize.curve\_fit}.
95\% bootstrap CIs on each $b$ from 10{,}000 resamples.  Confirmed
if (i)~the cross-profile slope range $b_{\max} - b_{\min}$ exceeds
$0.05$ per decade of corpus growth, and (ii)~Kendall's
$\tau < 1$ between the $5$K and $200$K profile orderings on
$\geq 1$ benchmark.  Pairwise slope differences tested via
permutation test (10{,}000 shuffles, $\alpha = 0.05$).

\paragraph{Power.}
For $n = 500$ queries with $\alpha_\text{corrected} = 0.00179$,
minimum detectable effect at 80\% power:
$d \approx 0.168$ (Cohen's $d$).

% ════════════════════════════════════════════════════════════════
\section{Exclusion Criteria}
\label{sec:exclusion}
% ════════════════════════════════════════════════════════════════

\begin{enumerate}[nosep,leftmargin=1.5em]
  \item Queries exceeding 30\,s wall-clock are excluded and counted.
  \item Configurations with $>$10\% null-retrieval queries are marked
    \textsc{Failed} and excluded from the ANOVA.
  \item Corpus sizes where index construction fails (OOM or $>$24\,h)
    are excluded from the scaling curve fit.
  \item Configurations require $\geq$200 valid queries (of 2{,}300)
    to enter the ANOVA.
\end{enumerate}

% ════════════════════════════════════════════════════════════════
\section{Budget, Timeline, and Deviations Protocol}
\label{sec:budget}
% ════════════════════════════════════════════════════════════════

\begin{table}[h]
\centering
\footnotesize
\caption{Estimated compute budget.}
\label{tab:budget}
\begin{tabular}{@{}lrrr@{}}
\toprule
Experiment & Configs & Queries/config & Est.\ cost (USD) \\
\midrule
Factorial sweep     & 242 & 2{,}300 & \$2{,}478 \\
Scaling experiment  & 15  & 500     & \$30 \\
Pareto experiment   & 20  & 500     & \$45 \\
\midrule
\textbf{Total}      &     &         & $\mathbf{\sim\$2{,}553}$ \\
\bottomrule
\end{tabular}
\end{table}

\textbf{Timeline.}
All experiments will be run \textit{after} archival of this document.
Week~1--3: factorial sweep (H1).
Week~4: scaling experiment (H3).
Week~4--5: Pareto experiment (H2).
Week~5--6: analysis and paper writing.

\textbf{Deviations protocol.}
Any deviation from this pre-registration will be: (1)~documented in a
dated amendment, (2)~reported transparently in the paper, and
(3)~accompanied by both the pre-registered analysis and the modified
analysis so that reviewers can assess the impact of the deviation.
Any analysis not specified here will be labeled \textit{exploratory}.

% ════════════════════════════════════════════════════════════════
\section*{Community Input}
% ════════════════════════════════════════════════════════════════

This pre-registration is published before experimental data collection
to establish hypothesis commitment and invite community input.
We welcome feedback on:
\begin{enumerate}[nosep,leftmargin=1.5em]
  \item additional hypotheses worth testing within this factorial
    framework,
  \item model selection for generation and evaluation
    (Section~\ref{sec:models}),
  \item benchmark composition and domain coverage,
  \item statistical methodology and effect-size thresholds.
\end{enumerate}
All community-suggested modifications adopted before experiments begin
will be documented as amendments; post-experiment suggestions will be
addressed as exploratory analyses in the companion paper.

% ════════════════════════════════════════════════════════════════
\bibliographystyle{plainnat}
\bibliography{preregistration}

\end{document}
